\documentclass[12pt]{article}

\usepackage{sbc-template}
\usepackage{graphicx,url}
\usepackage[utf8]{inputenc}
\usepackage[brazilian]{babel}
\usepackage{wasysym}
\usepackage[ruled,vlined]{algorithm2e}
\usepackage{fancyvrb}
\usepackage{float}
\usepackage{tikz}
\usetikzlibrary{positioning, calc, decorations.pathreplacing, arrows.meta}
\sloppy

\newcommand{\placeholderfigure}[3]{%
\begin{figure}[ht]
\centering
\fbox{%
\begin{minipage}[c][#1][c]{0.9\linewidth}
\centering
\textbf{#2}
\end{minipage}}
\caption{#3}
\end{figure}
}


\title{Exploiting Signature Fields to Fool PDF Signature Validation}

\author{Enzo R. Brum\inst{1}}


% TODO: FIX this thing
\address{ABCD -- EFGH\\
  Caixa Postal 15.064 -- 91.501-970 -- Porto Alegre -- RS -- Brazil
\nextinstitute
  LLLL\\
  LLLL
  %\email{\{nedel,flavio\}@inf.ufrgs.br, R.Bordini@durham.ac.uk,
  %jomi@inf.furb.br}
  \email{enzo.brum@ufsc.br}
}

\begin{document}

\maketitle

\begin{abstract}
	The Portable Document Format (PDF) is the standard way of sharing and defining electronic documents. It is a format used across a vast array of fields to represent official documents, many of which require strong guarantees that no unauthorized modifications are made to their contents. Digital signatures can be used to protect both the integrity and authenticity of a document, allowing an entity, whether a person, business, or service, to testify to one specific state of the document such that any non-permitted modification, for example altering the value of a contract or a health prescription, invalidates the signature. In this paper we introduce two novel attack classes that abuse flaws in the PDF specification to visually modify a page of a signed document without invalidating the signature or warning the user: Appearance Substitution Attack (ASA) and Signature Field Duplication Attack (SFDA). While previous works already explored incremental updates, we tap into modifications to the signature field itself, a specific kind of update that has yet to be fully explored. Our evaluation found that 8 out of 12 tested GUI applications and 2 popular verification libraries were vulnerable to at least one of the attack classes. We also propose countermeasures and specification improvements for fixing both.
\end{abstract}


\section{Introduction}

% TODO -> brainstorm about what made finding such error possible and use it to write intro -> justificativa

\begin{itemize}
	\item context for use of pdf signatures
	\item importance of them being considered invalid upon visual modification
	\item brief mention of incremental updates
	\item mention of attack breaking signatures
	\item main contributions
\end{itemize}

\begin{figure}[ht]
	\centering
	\begin{tikzpicture}[
			img/.style={inner sep=0pt, outer sep=0pt},
			lbl/.style={font=\small\itshape, yshift=-2pt},
		]
		\node[img] (og) {\includegraphics[width=0.47\linewidth]{og.png}};
		\node[img, right=0.6cm of og] (mod) {\includegraphics[width=0.47\linewidth]{mod.png}};
		\node[lbl, below=2pt of og]  {original signed contract};
		\node[lbl, below=2pt of mod] {after attack};
		\draw[->, very thick, >={Stealth[length=3mm, width=2mm]}]
		(og.east) -- (mod.west);
	\end{tikzpicture}
	\caption{Example of a contract signed on gov.br in which the contracted value has been altered from R\$\,100.000,00 to R\$\,200.000,00 after signing, while the signature is still reported as valid.\label{fig:attack-example}}
\end{figure}


\section{Background}

% TODO: may need to make it smaller if lacking space later on. good enough for now.

\subsection{Pdf Structure}

A PDF document consists of four parts: (i) a header declaring the PDF version; (ii) a body containing the document's objects; (iii) a cross-reference (xref) table mapping each object to its byte offset; and (iv) a trailer pointing to both the xref table and the document's root object.

The root object, the \texttt{Catalog}, is the entry point from which the entire object tree is traversed. It references a Pages dictionary, which links to the individual pages of the document, and an \texttt{AcroForm} dictionary, which describes any interactive form fields that users can fill or modify. An example of a basic PDF object tree can be seen on Figure~\ref{fig:pdf-object-tree}.


%\begin{noindent}
\begin{figure}[H]
	\centering
	\begin{Verbatim}[frame=single,fontsize=\footnotesize,xleftmargin=0.05\linewidth,xrightmargin=0.00\linewidth]
1 0 obj                 % object id 1
<<                      % start of dictionary
                        % entries have the form '/key value'
/Type     /Catalog      % Catalog object
/Pages    2 0 R         % ref to object 2 (Pages)
/AcroForm 3 0 R         % ref to object 3 (AcroForm)
>>                      % end of dictionary
endobj

2 0 obj  % Pages
<< /Kids [4 0 R]  /Count 1 >>
endobj

3 0 obj  % AcroForm
<< /Fields [6 0 R] >>
endobj

4 0 obj  % Page
<< /Parent 2 0 R  /Annots [6 0 R] >>
endobj
	\end{Verbatim}
	\caption{Skeleton of a PDF object tree.\label{fig:pdf-object-tree}}
\end{figure}
%\end{noindent}


\subsection{Incremental Updates}

When handling updates, a PDF processor can append changes to the end of the file rather than
overwriting existing content. As shown in Figure~\ref{fig:incremental-update}, modified objects are appended after the \texttt{\%EOF} marker, and a new xref table and trailer are added, creating a new revision of the document.

This mechanism preserves the original file content while incorporating updates. Most PDF readers merge all revisions into a unified view, though some (such as Adobe Acrobat
Reader~\cite{adobereader}) also expose tools for inspecting each revision individually.

%\begin{noindent}
\begin{figure}[ht]
    \centering
    \begin{tikzpicture}[
        box/.style={draw, minimum width=3cm, minimum height=0.7cm, align=center, font=\small},
        orig/.style={box, fill=blue!10},
        new/.style={box, fill=orange!30},
    ]
        % original revision (left column)
        \node[orig]                  (h1) {Header};
        \node[orig, below=0pt of h1] (b1) {Body};
        \node[orig, below=0pt of b1] (x1) {xref table};
        \node[orig, below=0pt of x1] (t1) {Trailer};

        % updated revision (right column)
        \node[orig, right=3.2cm of h1] (h2) {Header};
        \node[orig, below=0pt of h2]   (b2) {Body};
        \node[orig, below=0pt of b2]   (x2) {xref table};
        \node[orig, below=0pt of x2]   (t2) {Trailer};
        \node[new,  below=0pt of t2]   (bu) {Updated body};
        \node[new,  below=0pt of bu]   (xu) {New xref table};
        \node[new,  below=0pt of xu]   (tu) {New trailer};

        % arrow between columns, label above
        \coordinate (la) at ($(h1.east)!0.5!(t1.east)$);
        \coordinate (ra) at ($(h2.west)!0.5!(t2.west)$);
        \draw[->, thick] ($(la)+(0.25,0)$)
            -- node[above, font=\small\itshape] {incremental update}
            ($(ra)+(-0.25,0)$);
    \end{tikzpicture}
    \caption{Incremental update: the updated objects, a new xref table, and a new trailer are appended after the original revision.\label{fig:incremental-update}}
\end{figure}
%\end{noindent}


\subsection{Signature Field} \label{sec:sig-field}

A signature field is a specialized form field that acts as a container for a PDF digital signature. Each signature field is registered in the document's \texttt{AcroForm}dictionary and is associated with a single page. It holds a reference to the signature dictionary, which defines the digital signature itself, and can also specify how the signature is rendered visually on the page. This structure is illustrated in Figure~\ref{fig:sig-field-structure}.

%\begin{noindent}
\begin{figure}[H]
    \centering
    \begin{tikzpicture}[
        dict/.style={draw, rectangle, align=left, inner sep=5pt, minimum width=3cm},
        arr/.style={->, thick, >=stealth},
    ]
        % AcroForm dictionary
        \node[dict] (acro) {%
            \begin{tabular}{@{}l@{}}
                \textbf{\sffamily AcroForm}\\[2pt]
                \texttt{/Fields [...]}
            \end{tabular}
        };

        % Signature field
        \node[dict, right=1.0cm of acro] (field) {%
            \begin{tabular}{@{}l@{}}
                \textbf{\sffamily Signature field}\\[2pt]
                \texttt{/T \ \ \ Sig1}\\
                \texttt{/V \ \ \ ...}\\
                \texttt{/Rect [...]}\\
                \texttt{/AP \ \ ...}\\
                \texttt{/SigFieldLock ...}
            \end{tabular}
        };

        % Signature dictionary
        \node[dict, right=1.0cm of field] (sigdict) {%
            \begin{tabular}{@{}l@{}}
                \textbf{\sffamily Signature dictionary}\\[2pt]
                \texttt{/Type \ \ /Sig}\\
                \texttt{/ByteRange [...]}\\
                \texttt{/Contents ...}\\
            \end{tabular}
        };

        % references
        \draw[arr] (acro.east)  -- (field.west);
        \draw[arr] (field.east) -- (sigdict.west);
    \end{tikzpicture}
    \caption{Signature-field structure: the \texttt{AcroForm}registers each signature field, which references its signature dictionary via \texttt{/V}.\label{fig:sig-field-structure}}
\end{figure}
%\end{noindent}

Its most relevant attributes are \texttt{/T}, a name that uniquely identifies the field within the document; \texttt{/V}, a reference to the signature dictionary; \texttt{/AP}, which defines the field's visual appearance; \texttt{/Rect}, which specifies the field's position and dimensions on the page; and \texttt{/SigFieldLock}, which is discussed in Section~\ref{sec:doc-mod-perm}.

Setting the \texttt{/V} entry of an existing signature field is defined by the PDF specification as \emph{filling} that field.

\subsection{Digital Signatures}

A document is digitally signed by appending a new signature field, linked to a page and to a new
signature dictionary, as an incremental update. Each signature dictionary contains the user's
certificate and a byte range covering the full document as it existed at signing time; the digital signature is computed over that byte range using the user's private key and stored within the signature dictionary.

This design naturally supports multiple signatures. As shown in Figure~\ref{fig:multiple-signatures}, a subsequent signature can be appended via another incremental update without invalidating any prior signature since the bytes covered by those earlier signatures remain unchanged.

%\begin{noindent}
\begin{figure}[H]
    \centering
    \begin{tikzpicture}[
        >={Stealth[length=2mm, width=1.6mm]},
        box/.style={draw, line width=0.4pt, rounded corners=1pt,
                    minimum width=2.2cm, minimum height=0.55cm, align=center, font=\small},
        orig/.style={box, fill=blue!10},
        sig1/.style={box, fill=orange!30},
        sig2/.style={box, fill=green!25},
        brc1/.style={decorate, decoration={brace, amplitude=4pt, raise=2pt, mirror},
                     thick, draw=orange!70!black},
        brc2/.style={decorate, decoration={brace, amplitude=4pt, raise=2pt, mirror},
                     thick, draw=green!50!black},
        brlbl1/.style={font=\footnotesize, text=orange!70!black},
        brlbl2/.style={font=\footnotesize, text=green!50!black},
    ]
        % col 1: unsigned document
        \node[orig]                  (Uh) {Header};
        \node[orig, below=0pt of Uh] (Ub) {Body};
        \node[orig, below=0pt of Ub] (Ux) {xref};
        \node[orig, below=0pt of Ux] (Ut) {Trailer};

        % col 2: state after signing with Sig 1
        \node[orig, right=2.7cm of Uh]  (Ah)  {Header};
        \node[orig, below=0pt of Ah]    (Ab)  {Body};
        \node[orig, below=0pt of Ab]    (Ax)  {xref};
        \node[orig, below=0pt of Ax]    (At)  {Trailer};
        \node[sig1, below=0pt of At]    (As)  {Signature 1};
        \node[sig1, below=0pt of As]    (Axn) {new xref};
        \node[sig1, below=0pt of Axn]   (Atn) {new trailer};

        % col 3: state after signing with Sig 2
        \node[orig, right=2.7cm of Ah]  (Bh)   {Header};
        \node[orig, below=0pt of Bh]    (Bb)   {Body};
        \node[orig, below=0pt of Bb]    (Bx)   {xref};
        \node[orig, below=0pt of Bx]    (Bt)   {Trailer};
        \node[sig1, below=0pt of Bt]    (Bs1)  {Signature 1};
        \node[sig1, below=0pt of Bs1]   (Bxn1) {new xref};
        \node[sig1, below=0pt of Bxn1]  (Btn1) {new trailer};
        \node[sig2, below=0pt of Btn1]  (Bs2)  {Signature 2};
        \node[sig2, below=0pt of Bs2]   (Bxn2) {new xref};
        \node[sig2, below=0pt of Bxn2]  (Btn2) {new trailer};

        % col 2 brace: Sig 1 byte range
        \draw[brc1] ([xshift=4pt]Atn.south east) -- ([xshift=4pt]Ah.north east)
            node[midway, right=8pt, brlbl1] {$Sig1$ range};

        % col 3 braces: inner Sig 1, outer Sig 2
        \draw[brc1] ([xshift=4pt]Btn1.south east)  -- ([xshift=4pt]Bh.north east)
            node[midway, right=8pt, brlbl1] {$Sig1$ range};
        \draw[brc2] ([xshift=60pt]Btn2.south east) -- ([xshift=58pt]Bh.north east)
            node[midway, right=8pt, brlbl2] {$Sig2$ range};

        \draw[->, thick] ([yshift=-30pt]Uh.north east) -- ([yshift=-30pt]Ah.north west)
            node[midway, above, font=\footnotesize\itshape] {$Sig1$};
        \draw[->, thick] ([yshift=-30pt]Ah.north east) -- ([yshift=-30pt]Bh.north west)
            node[midway, above, font=\footnotesize\itshape] {$Sig2$};
    \end{tikzpicture}
    \caption{Each signing applies an incremental update; the new signature covers all bytes that existed at signing time, so earlier signatures remain valid.\label{fig:multiple-signatures}}
\end{figure}
%\end{noindent}

\subsection{Modification Permission Mechanism} \label{sec:doc-mod-perm}

% TODO: read ISO again just to be sure. my curr understanding is that spec does not define incr updates verification and that what readers do isn't required.

Because incremental updates preserve all prior bytes, visual content can be modified after signing without technically invalidating the signature. An attacker could, for example, alter a page or inject an annotation via an incremental update while leaving the signed byte range intact, all within the boundaries of the specification~\cite{iso32000-2}. In response, most verifiers~\cite{adobereader,foxitreader,libreoffice} either treat such post-signing modifications as invalidating the signature or surface a warning to the user.

The specification also provides mechanisms that make signatures sensitive to specific classes of modifications. \texttt{DocMDP} defines which types of modification are permitted before the signature is considered broken; \texttt{FieldMDP} locks specific form fields so that any change to them invalidates the signature; and \texttt{SigFieldLock} can do both capabilities. The differences are summarized in Figure~\ref{fig:mdp-mechanisms}.

\begin{figure}[ht]
	\centering
	% (a) capabilities
	\begin{minipage}{\linewidth}
		\centering
		\small
		\begin{tabular}{@{}l p{4.8cm} l@{}}
			\hline
			\textbf{Mechanism}    & \textbf{Scope}                                          & \textbf{Allowed in}           \\
			\hline
			\texttt{DocMDP}       & whole document, per permission level                    & certification signatures only \\
			\texttt{FieldMDP}     & listed form fields only                                 & any signature                 \\
			\texttt{SigFieldLock} & listed form fields, plus whole document (since PDF~2.0) & any signature                 \\
			\hline
		\end{tabular}\\[3pt]
		{\footnotesize (a) Capabilities of each mechanism.}
	\end{minipage}\\[10pt]
	% (b) permission levels
	\begin{minipage}{\linewidth}
		\centering
		\small
		\begin{tabular}{@{}c l@{}}
			\hline
			\textbf{Level} & \textbf{Permitted modifications}                              \\
			\hline
			1              & none                                                          \\
			2              & form fill-in and adding further signatures                    \\
			3              & level 2, plus annotation creation, modification, and deletion \\
			\hline
		\end{tabular}\\[3pt]
		{\footnotesize (b) Permission levels, shared by \texttt{DocMDP} and \texttt{SigFieldLock}.}
	\end{minipage}
	\caption{PDF modification-permission mechanisms: (a) what each mechanism protects and which signatures may carry it; (b) the permission-level legend shared by \texttt{DocMDP} and \texttt{SigFieldLock}.\label{fig:mdp-mechanisms}}
\end{figure}

\texttt{DocMDP} and \texttt{FieldMDP} reside in the signature dictionary; \texttt{SigFieldLock} is referenced by the signature field itself, and its value is copied into \texttt{FieldMDP} at signing time.

As described in~\cite{pdf-certification-attack}, PDF distinguishes between \emph{approval} signatures, which attest to the document's state, and a \emph{certification} signature, which additionally governs which modifications are permissible. A document may carry at most one certification signature, which must appear first; any number of approval signatures may follow. \texttt{DocMDP} is exclusive to certification signatures. \texttt{FieldMDP} and \texttt{SigFieldLock} may be used by either kind. Since PDF 2.0~\cite{iso32000-2}, \texttt{SigFieldLock} carries the same expressive power as \texttt{DocMDP}, allowing approval signatures to impose document-level modification restrictions as well.


\section{Threat Model}

The attacker's goal is to alter the visual content of a signed PDF while ensuring that signature verifiers continue to report the original signature as valid and raise no warning about post-signing modifications. We assume the attacker has obtained a PDF carrying a digital signature that the target verifier already trusts and considers valid.

The \emph{victim} is the person, organization, or service that needs to establish whether a document carries a trusted signature. The \emph{signer} is the party that originally signed the document and, in doing so, configured any modification-permission mechanisms.

We consider two classes of verifier: graphical user interface (GUI) applications and programming libraries. The distinction matters because each class exposes a different surface to the user, and our attacks succeed in different ways on each.

\subsection{User Interface Verifiers}

We adopt a layered model for GUI verifiers similar to that of~\cite{pdf-certification-attack}, but omitting the PDF-annotation layer since our attacks do not exercise it.

\textbf{UI-Layer 1} is the top-bar validation status. Most GUI verifiers display this banner immediately upon opening the document and use it to summarize whether all signatures in the document are valid. Owing to its prominence, it is typically the first cue a user inspects when assessing the document's authenticity.

\textbf{UI-Layer 2} carries detailed per-signature information, such as the signer's name and certificate chain. It is generally not surfaced automatically; the user must open a dedicated panel before this information becomes visible.

\subsection{Programming Libraries}

Programming libraries expose a binary verdict and a detailed report describing every check performed for each verified signature. We collapse both into a single abstraction, the \textbf{API-Layer}, since they are returned together in the same function response and any developer consuming the library has access to both at once.

In this abstraction, the API-Layer answers two questions: whether the signature is considered valid, and whether any warning should be noted alongside the verdict. For instance, one of such warnings can be that the verifier detected an uncommon or post-signing modification but still accepts the signature as valid.

\subsection{Winning Conditions} \label{sec:winning-conditions}

For each UI layer for GUI verifiers or the single API-Layer for library verifiers, we classify the outcome of an attack as follows. The output is \emph{vulnerable} (\CIRCLE) when it reports the signature as valid without flagging a post-signing modification. It is \emph{partially vulnerable} (\LEFTcircle) when it reports the signature as valid but raises a warning that the document was modified. It is \emph{secure} (\Circle) when it flags the signature as invalid.

For GUI verifiers, an attack is \emph{successful} when both UI layers are vulnerable, and \emph{partially successful} when UI-Layer~1 is vulnerable while UI-Layer~2 is either \emph{partially vulnerable} or \emph{secure}. We treat partial success as meaningful because an inattentive user may rely solely on the top-bar banner and overlook the warning surfaced in the detail panel, accepting a modified document despite the recorded discrepancy. For library verifiers, an attack is \emph{successful} when the API-Layer is vulnerable, and \emph{partially successful} when it is partially vulnerable; the same reasoning applies, since a developer may consume only the binary verdict and ignore the accompanying warnings.


\section{Breaking Digital Signature Verifiers} \label{sec:attack-defintions}

Our attacks exploits two underspecified aspects of the PDF specification~\cite{iso32000-2} to alter existing signature fields and use them to introduce visual changes. By modifying a signature field's existing entries, an attacker can place arbitrary visual content on a page. Consequently, attacks such as the one shown in Figure~\ref{fig:attack-example} can be carried out without invalidating the signature in affected verifiers.

% TODO: cant use ASA here maybe? pdf-insecurity paper on certification already defines sneaky signature attack
\subsection{Appearance Substitution Attack (ASA)}

% 1. reexplain /Rect and /AP
% 2. explain that /Rect can be resized and /AP can be overwritten as to change what's actualy shown on page. -> point to figure with an already signed document with a common stamp -> change it so rect compasses whole page, the stamp is kept and text is changed. show virtual square outline so people see it
% 3. argument the spec is flawed, provide exact wording and say it does NOT necessarily restrict such attack.

% TODO: fact check invisible signatures.

As described in Section~\ref{sec:sig-field}, a signature field's \texttt{/Rect} entry defines the region of the page occupied by the field, while its \texttt{/AP} entry defines the appearance drawn inside that region. For signatures with a visible component, \texttt{/Rect} typically identifies a small area of the page containing the signer's name or stamp, and \texttt{/AP} stores the corresponding visual representation. For signatures without a visible component, \texttt{/Rect} commonly contains a zero-area rectangle, and \texttt{/AP} is absent.

These entries can be abused to turn a signature field into a vehicle for arbitrary visual content. Expanding \texttt{/Rect} to cover a larger region, or even the entire page, increases the area in which the field appearance can be rendered. Replacing \texttt{/AP} then changes what is drawn in that region. As shown in Figure~\ref{fig:rect-ap-abuse}, an attacker can start from a document that already contains a visible signature, preserve the original signature appearance, and add a small malicious modification elsewhere on the page by carefully crafting the replacement appearance stream.

As discussed in Section~\ref{sec:doc-mod-perm}, permission levels two and three of \texttt{DocMDP} and \texttt{SigFieldLock} allow form fields to be filled through an incremental update without invalidating the protected signature. Since a signature field is itself a form field, the specification does not explicitly rule out filling such a field after a prior signature has been applied. The relevant ISO wording further leaves room for this interpretation:

\begin{quote}
	\emph{[Filling in (signing) the signature field entails updating at least the \texttt{/V} entry and usually also the \texttt{/AP} entry of the associated widget annotation. Exporting a signature field typically exports the \texttt{/T}, \texttt{/V}, and \texttt{/AP} entries. (Section 12.7.5.5 of ISO~32000-2~\cite{iso32000-2})]}
\end{quote}

The phrase ``at least'' is important: it suggests that filling a signature field is not necessarily limited to updating \texttt{/V} and \texttt{/AP}. Under a permissive interpretation, other field entries, including \texttt{/Rect}, could also be treated as part of the same filling operation. This ambiguity enables ASA to be applied to documents protected by modification-permission mechanisms while remaining within the class of changes that affected verifiers accept as non-invalidating.

\begin{figure}[H]
	\centering
	% (a) visual outcome
	\begin{minipage}{\linewidth}
		\centering
		\begin{tikzpicture}[
				img/.style={inner sep=0pt, outer sep=0pt},
				lbl/.style={font=\small\itshape, yshift=-2pt},
			]
			\node[img] (og)  {\includegraphics[width=0.44\linewidth]{og_outlined_stamp_rect.png}};
			\node[img, right=0.6cm of og] (mod) {\includegraphics[width=0.44\linewidth]{mod_outlined.png}};
			\node[lbl, below=2pt of og]  {original signed contract};
			\node[lbl, below=2pt of mod] {after ASA};
			\draw[->, very thick, >={Stealth[length=3mm, width=2mm]}]
			(og.east) -- (mod.west);
		\end{tikzpicture}\\[3pt]
		{\footnotesize (a) Visual outcome. Cyan outlines mark the signature field's \texttt{/Rect}: a small region around the original stamp on the left, expanded to cover the whole page on the right. Red outlines on the right highlight the original stamp redrawn at its original location and the malicious value that overwrites \texttt{R\$\,100.000,00} with \texttt{R\$\,200.000,00}.}
	\end{minipage}\\[12pt]
	% (b) structural view
	\begin{minipage}{\linewidth}
		\centering
		\begin{tikzpicture}[
				dict/.style={draw, rectangle, align=left, inner sep=5pt, font=\footnotesize},
				note/.style={font=\small\itshape},
			]
			% original field
			\node[dict] (ogfield) {%
				\begin{tabular}{@{}l@{}}
					\textbf{\sffamily Signature field} \\[2pt]
					\texttt{/T \ \ \ Sig1}             \\
					\texttt{/V \ \ \ <sig dict>}       \\
					\texttt{/Rect [small area]}        \\
					\texttt{/AP \ \ <original AP, draws stamp>}
				\end{tabular}
			};

			% modified field
			\node[dict, right=2.0cm of ogfield] (modfield) {%
				\begin{tabular}{@{}l@{}}
					\textbf{\sffamily Signature field (after ASA)} \\[2pt]
					\texttt{/T \ \ \ Sig1}                         \\
					\texttt{/V \ \ \ <sig dict>}                   \\
					\texttt{/Rect [whole page]}                    \\
					\texttt{/AP \ \ <new AP, redraws stamp +}      \\
					\hspace*{5em}\texttt{overwrites value>}
				\end{tabular}
			};

			\draw[->, very thick, >={Stealth[length=3mm, width=2mm]}]
			(ogfield.east) -- node[above, note] {ASA} (modfield.west);
		\end{tikzpicture}\\[3pt]
		{\footnotesize (b) Abstract structural change. \texttt{/T} and \texttt{/V} are preserved, while \texttt{/Rect} is enlarged to span the whole page and a new \texttt{/AP} is created that redraws the original stamp at its original location and overwrites the contracted value elsewhere on the page.}
	\end{minipage}
	\caption{ASA applied to a contract signed at gov.br: (a) visual example of the attack; (b) abstract structural view of the modified signature field.\label{fig:rect-ap-abuse}}
\end{figure}


If the existing signature field has no visual component, a new one can easily be added by creating a new \texttt{/AP} and modifying \texttt{/Rect}.

\subsection{Signature Field Duplication Attack (SFDA)}

A related attack arises from the opposite direction: instead of mutating an existing signature field, the attacker introduces a \emph{new} signature field through an incremental update and sets its \texttt{/V} entry to point at the signature dictionary of a signature already present in the document. The duplicated field is registered in the \texttt{AcroForm}dictionary alongside the original and can be placed on any page, with attacker-controlled \texttt{/Rect} and \texttt{/AP} entries. Because the bytes covered by the original signature are untouched, the cryptographic verification of that signature continues to succeed, and affected verifiers render the duplicated field as a second valid signature. An example of the resulting document is shown in Figure~\ref{fig:sfda-example}.

SFDA is closely related to the Sneaky Signature Attack of~\cite{pdf-certification-attack} but differs in one fundamental respect: it does not require the attacker to produce a fresh signature trusted by the verifier. Whereas the Sneaky Signature Attack appends an attacker-controlled signature dictionary, SFDA reuses the legitimate one. This makes the attack particularly damaging in scenarios where obtaining a trusted certificate is impractical, such as documents signed under closed PKIs or restricted issuance policies.

Because the duplicated field inherits the certificate and signer name of the original, the forgery is harder for a human reviewer to dismiss than one bearing an unfamiliar identity. The duplication itself is visually detectable, since two signers now appear in the verifier interface, but the verifier displays a known name and a familiar certificate chain, which is precisely what an inattentive reviewer is most likely to accept.

The PDF specification does not prohibit multiple signature fields from referencing the same signature dictionary. ISO~32000-2 acknowledges the configuration but treats it only as a matter of presentation:

\begin{quote}
	\emph{[If more than one location is associated with a signature, the meaning can become ambiguous. (Section 12.7.5.5 of ISO~32000-2~\cite{iso32000-2})]}
\end{quote}


% TODO: add figure
% TODO: not exactly true, go back later.
The text stops short of disallowing the configuration or prescribing how verifiers should handle it. Affected verifiers resolve the ambiguity by rendering every signature field that references the signature dictionary as valid, the attacker-controlled one included, while still reporting the document as correctly signed.

\begin{figure}[ht]
	\centering
	% (a) visual outcome
	\begin{minipage}{\linewidth}
		\centering
		\begin{tikzpicture}[
				img/.style={inner sep=0pt, outer sep=0pt},
				lbl/.style={font=\small\itshape, yshift=-2pt},
			]
			\node[img] (og)  {\includegraphics[width=0.44\linewidth]{layer_2.png}};
			\node[img, right=0.6cm of og] (mod) {\includegraphics[width=0.44\linewidth]{mod_layer_2.png}};
			\node[lbl, below=2pt of og]  {original verifier output};
			\node[lbl, below=2pt of mod] {after SFDA};
			\draw[->, very thick, >={Stealth[length=3mm, width=2mm]}]
			(og.east) -- (mod.west);
		\end{tikzpicture}\\[3pt]
		{\footnotesize (a) Visual outcome at Layer-2. The verifier on the left correctly reports two valid signatures on a document carrying two signatures. After SFDA (right), the same verifier reports three valid signatures: the second signature has been duplicated, and the verifier accepts the duplicate as a third independent valid signature.}
	\end{minipage}\\[12pt]
	% (b) structural view
	\begin{minipage}{\linewidth}
		\centering
		\begin{tikzpicture}[
				box/.style={draw, rectangle, align=center, inner sep=4pt, font=\footnotesize, minimum width=1.8cm},
				newbox/.style={box, fill=orange!20, draw=orange!70!black, thick},
				arr/.style={->, thick, >=stealth},
				newarr/.style={->, thick, >=stealth, draw=orange!70!black},
			]
			% AcroForm
			\node[box, minimum width=8cm] (acro) {%
			\textbf{\sffamily AcroForm}\quad\texttt{/Fields [Sig1, Sig2,} \textcolor{orange!70!black}{\texttt{Sig2\_dup}}\texttt{]}%
			};

			% Fields
			\node[box, below=0.7cm of acro, xshift=-2.8cm] (s1) {Sig1 field};
			\node[box, below=0.7cm of acro]                (s2) {Sig2 field};
			\node[newbox, below=0.7cm of acro, xshift=2.8cm] (sdup) {%
				\begin{tabular}{@{}c@{}}
					Sig2\_dup field                    \\[1pt]
					\emph{\scriptsize attacker-chosen} \\
					\emph{\scriptsize \texttt{/Rect}, \texttt{/AP}}
				\end{tabular}
			};

			% Dicts
			\node[box, below=1.2cm of s1] (d1) {Dict1};
			\node[box, below=1.2cm of s2] (d2) {Dict2};

			% AcroForm -> fields
			\draw[arr]    (acro.south -| s1.north)   -- (s1.north);
			\draw[arr]    (acro.south -| s2.north)   -- (s2.north);
			\draw[newarr] (acro.south -| sdup.north) -- (sdup.north);

			% Fields -> dicts
			\draw[arr] (s1) -- node[right, font=\scriptsize] {\texttt{/V}} (d1);
			\draw[arr] (s2) -- node[right, font=\scriptsize] {\texttt{/V}} (d2);
			\draw[newarr] (sdup.south) to[out=-90, in=0]
			node[above right, font=\scriptsize] {\texttt{/V}} (d2.east);
		\end{tikzpicture}\\[3pt]
		{\footnotesize (b) Abstract structural change (new field and edges in orange). A new signature field \texttt{Sig2\_dup} is appended to the AcroForm's \texttt{/Fields} array; its \texttt{/V} entry points at \texttt{Dict2}, the same signature dictionary already referenced by \texttt{Sig2}, while its \texttt{/Rect} and \texttt{/AP} are chosen by the attacker. The field is registered on an attacker-chosen page via that page's \texttt{/Annots} array.}
	\end{minipage}
	\caption{SFDA applied to a document containing two valid signatures: (a) visual example of the attack; (b) abstract structural view of the duplicated signature field.\label{fig:sfda-example}}
\end{figure}

\subsection{Limitations}

ASA can only modify an existing signature field, and is therefore constrained by the page on which that field is presented. Documents that place signature fields on a dedicated signature page consequently offer a smaller attack surface than documents where signature fields appear alongside important content. SFDA is not bound by this same placement constraint, but it introduces a duplicated signature field that remains visible within UI-layer~2.

Given both attacks abuse visual components of interactive fields, their rendering can vary across PDF readers. Some readers may highlight the field, make it clickable, or prevent text selection inside it. These artifacts can make the attack detectable by an attentive user, but prior work shows that users may still perform security-sensitive actions, such as entering their own credentials, when the inconsistencies are small~\cite{phishing-password-managers}. Similarly, in our setting, users may interpret these artifacts as quirks of the reader or verifier, especially when all signatures are reported as valid. This is also similar to prior attacks on PDF signatures, where the abused functionality left subtle visual inconsistencies~\cite{pdf-certification-attack}.



\section{Methodology}

We evaluate four GUI verifiers (Adobe Acrobat~\cite{adobereader}, Foxit PDF Reader~\cite{foxitreader}, Master PDF Editor~\cite{masterpdfeditor}, and Microsoft Edge~\cite{microsoftedge}), two library verifiers (pyHanko~\cite{pyhanko} and the European Commission's Digital Signature Services~\cite{eudss}), and eight web verifiers (the DSS Demonstration WebApp~\cite{dssdemo}, the ITI signature validator~\cite{itivalidador}, Certisign~\cite{certisign}, DigitalSign~\cite{digitalsign}, the Finnish DVV validator~\cite{dvv}, BRY Verifica~\cite{bry}, ZapSign~\cite{zapsign}, and Clicksign~\cite{clicksign}) against the Appearance Substitution Attack (ASA) and the Signature Field Duplication Attack (SFDA), classifying each outcome under the criteria of Section~\ref{sec:winning-conditions}.

We use two different testing regimes. For library verifiers and GUI desktop verifiers, where the trust anchors can be freely chosen, an automated tester (Section~\ref{sec:automated-tester}) evaluates every verifier against many possible attack configurations. For GUI web verifiers, whose trusted-certificate lists are fixed by the creator and signatures with arbitrary parameters cannot be easily created, we perform a manual evaluation in which a single document carrying a signature accepted by the verifier is subjected to each attack and the outcome recorded.

For each verifier and attack, we record the UI-Layer~1 and UI-Layer~2 outcomes (for GUI verifiers) or the API-Layer outcome (for library verifiers) and label each as \emph{vulnerable}, \emph{partially vulnerable}, or \emph{secure} following Section~\ref{sec:winning-conditions}.


\subsection{Automated Tester} \label{sec:automated-tester}

The automated tester is composed of two execution flows: test-case generation and test execution. Test-case generation enumerates configurations of the trusted signature and, for each configuration, produces one PDF carrying the corresponding ASA attack and one PDF carrying the corresponding SFDA attack. Each generated file is persisted, together with its expected veredict, in a local database for use by later runs.

To generate test cases, we enumerate all combinations of the following parameters, except for the invalid case in which the signature is a certification signature and the MDP level is not defined.

\begin{itemize}
	\item $is\_certification\_signature \in \{true, false\}$: defines whether the trusted signature is a \texttt{certification} signature or an \texttt{approval} signature.
	\item $mdp\_level \in \{null, 1, 2, 3\}$: defines the permission level. If it is \emph{null} and the signature is an approval signature, no protection mechanism is included. If it is not \emph{null}, certification signatures use \texttt{/DocMDP}, while approval signatures use \texttt{SigFieldLock}, both configured with the selected permission level.
	\item $is\_signature\_field\_protected \in \{true, false\}$: defines whether the trusted signature field is protected by \texttt{FieldMDP}.
	\item $is\_stamp\_enabled \in \{true, false\}$: defines whether the trusted signature has a preexisting visual component.
	\item $change\_page \in \{true, false\}$: defines whether the modified or duplicated signature field should be placed on the same page as the original field.
\end{itemize}

Test execution differs by verifier class. For library verifiers, each test case is submitted through the relevant verification function and the returned veredict and modification warnings are stored directly, populating the API-Layer from the layered model. For GUI desktop verifiers, the tester follows the screenshot-based approach of previous works~\cite{pdf-certification-attack}: it opens the document in the verifier, captures one screenshot per UI layer being evaluated, and classifies each by template matching against a small reference-image set per verifier and layer, one set per outcome class (\emph{valid}, \emph{warn}, \emph{invalid}). The \emph{warn} class captures the partial-vulnerability case of Section~\ref{sec:winning-conditions}, in which the verifier accepts the signature but raises a warning that the document was modified after signing.

\section{Evaluation} \label{sec:evaluation}
\begin{table}[ht]
	\centering
	\caption{Per-verifier results of how each verifier handles ASA and SFDA. \CIRCLE means a verifier is fully \emph{vulnerable} and doesn't warn the user that there were modifications to the document despite the signature being valid. \LEFTcircle means a verifier is \emph{partially vulnerable} and warns the user that a modification occured. \Circle means a verifier is \emph{secure} and correctly invalidates the signature. Empty cells mark layers that do not apply to the verifier's class or versions not exposed by the verifier.}
	\footnotesize
	\begin{tabular}{@{}lllccc@{}}
		\hline
		\textbf{Application}                                        & \textbf{Version} & \textbf{Category} & \textbf{UI-Layer 1}            & \textbf{UI-Layer 2}            & \textbf{API-Layer}         \\
		                                                            &                  &                   & ASA\hspace{4pt}SFDA            & ASA\hspace{4pt}SFDA            & ASA\hspace{4pt}SFDA        \\
		\hline
		Adobe Acrobat                                               & 26.1.21563.0     & desktop           & \LEFTcircle\hspace{4pt}\CIRCLE & \LEFTcircle\hspace{4pt}\CIRCLE &                            \\
		Foxit PDF Reader                                            & 2026.1.1.36485   & desktop           & \CIRCLE\hspace{4pt}\CIRCLE     & \LEFTcircle\hspace{4pt}\Circle &                            \\
		Master PDF Editor                                           & 5.9.9.8          & desktop           & \CIRCLE\hspace{4pt}\CIRCLE     & \CIRCLE\hspace{4pt}\CIRCLE     &                            \\
		Microsoft Edge                                              & 148.0.3967.83    & desktop           & \CIRCLE\hspace{4pt}\CIRCLE     & \CIRCLE\hspace{4pt}\CIRCLE     &                            \\
		pyHanko                                                     & 0.31.0           & library           &                                &                                & \Circle\hspace{4pt}\CIRCLE \\
		DSS                                                         & 6.4              & library           &                                &                                & \CIRCLE\hspace{4pt}\Circle \\
		DSS Demo WebApp                                             & 6.4              & web               & \LEFTcircle\hspace{4pt}\Circle & \LEFTcircle\hspace{4pt}\Circle &                            \\
		ITI Validador                                               & 6aec769          & web               & \CIRCLE\hspace{4pt}\Circle     & \CIRCLE\hspace{4pt}\Circle     &                            \\
		Certisign                                                   &                  & web               & \CIRCLE\hspace{4pt}\CIRCLE     & \CIRCLE\hspace{4pt}\CIRCLE     &                            \\
		DigitalSign                                                 &                  & web               & \CIRCLE\hspace{4pt}\Circle     & \CIRCLE\hspace{4pt}\Circle     &                            \\
		DVV                                                         &                  & web               & \LEFTcircle\hspace{4pt}\Circle & \LEFTcircle\hspace{4pt}\Circle &                            \\
		BRY Verifica                                                &                  & web               & \CIRCLE\hspace{4pt}\LEFTcircle & \CIRCLE\hspace{4pt}\LEFTcircle &                            \\
		ZapSign                                                     &                  & web               & \Circle\hspace{4pt}\Circle     & \Circle\hspace{4pt}\Circle     &                            \\
		Clicksign                                                   &                  & web               & \CIRCLE\hspace{4pt}\CIRCLE     & \CIRCLE\hspace{4pt}\CIRCLE     &                            \\
		\hline
		\multicolumn{3}{@{}l}{\textbf{ASA, vulnerable or partial}}  & 11/12            & 11/12             & 1/2                                                                                          \\
		\multicolumn{3}{@{}l}{\textbf{SFDA, vulnerable or partial}} & 7/12             & 6/12              & 1/2                                                                                          \\
		\hline
	\end{tabular}
	\label{tab:results}
\end{table}


This section presents the results of our analysis, covering 8 web-based, 4 desktop, and 2 library verifiers. Desktop and library verifiers, which support automation, were each evaluated against 112 crafted PDF documents; web verifiers were evaluated manually with a single document carrying a signature trusted by each verifier. Desktop tests ran on Windows 11, web tests on Windows 11 with Microsoft Edge, and library tests on NixOS 26. Signatures for the automated tests were produced by a custom pyHanko script using a self-issued certificate. The results of the evaluation can be seen at Table~\ref{tab:results}.

We first analyze how ASA performs against the evaluated verifiers. As Table~\ref{tab:results} shows, the large majority of verifiers are vulnerable to ASA at \textbf{UI-Layer~1} or the \textbf{API-Layer}, the layers users rely on to judge whether a signature is valid. With the sole exception of Foxit~\cite{foxitreader}, every GUI verifier vulnerable to ASA at \textbf{UI-Layer~1} is also vulnerable at \textbf{UI-Layer~2}, so the detail panel rarely catches the attack even when the user opens it. By the criteria of Section~\ref{sec:winning-conditions}, this yields 8 \emph{successful} attacks and 4 \emph{partially successful} ones.

We then analyze how SFDA performs against the evaluated verifiers. SFDA reaches slightly fewer verifiers than ASA, with exactly half of the GUI and library verifiers reporting the signature as valid in some form at \textbf{UI-Layer~1} or the \textbf{API-Layer}. As with ASA, Foxit is the only GUI verifier whose veredict downgrades between UI layers, going from fully vulnerable at \textbf{UI-Layer~1} to \emph{secure} at \textbf{UI-Layer~2}. Overall, SFDA yields 6 \emph{successful} attacks and 2 \emph{partially successful} ones.

Foxit's downgrade between UI layers reflects how it renders \textbf{UI-Layer~1}. Inspecting the automated test outputs, we observed that even when the underlying checks detect post-signing modifications, the top bar shows only a green ``certified document'' icon and no further detail. The warnings driving the partial \textbf{UI-Layer~2} veredict for ASA, and the \emph{secure} veredict for SFDA, only appear once the user opens the per-signature detail panel.

The page-change variants of the two attacks behave asymmetrically. Every verifier vulnerable to SFDA at \textbf{UI-Layer~1} remains vulnerable when the duplicated signature field is placed on a different page. This is arguably the stronger variant, since it lets the attacker stage the forged appearance over arbitrary content on any page. The same does not hold for ASA, where DSS~\cite{eudss}, ITI~\cite{itivalidador}, BRY~\cite{bry}, and DigitalSign~\cite{digitalsign} all reject the attack once the modified signature field is relocated to a different page.

pyHanko~\cite{pyhanko} is vulnerable to SFDA only when neither DocMDP nor SigFieldLock is present, and rejects both attacks otherwise. DSS exhibits a similar behavior for ASA, accepting the attack only when the modified signature field originally defined an invisible signature. Other verifiers built on DSS, such as the DSS Demo WebApp~\cite{dssdemo} and DVV~\cite{dvv} among those we tested, are plausibly vulnerable in the same configuration, but we could not confirm it, as no documents carrying an invisible signature trusted by these verifiers were available. ZapSign~\cite{zapsign} rejects both attacks, but does so by treating any post-signing modification as invalidating, including benign ones. In our tests, even appending a comment or a new revision to a signed document was enough to invalidate the signature.








\section{Countermeasures}

We propose countermeasures for both attacks at three levels: the specification, individual verifiers, and external detection tooling.

As outlined in Section~\ref{sec:attack-defintions}, both attacks exploit ambiguities in the PDF specification that allow a document's visual content to be modified without invalidating the protecting signature. ASA leverages the absence of a precise definition of what \emph{filling} a signature field entails, which lets verifiers treat post-signing changes to \texttt{/AP} and \texttt{/Rect} as non-invalidating. SFDA leverages the absence of any prescribed behaviour when multiple signature fields reference the same signature dictionary, which lets an attacker introduce a duplicate field that the verifier renders as a valid signature.

The specification can close both gaps. To prevent ASA, ISO~32000-2 should require that any modification to the \texttt{/AP} or \texttt{/Rect} of a filled signature field invalidates the protecting signature, regardless of the prevailing \texttt{DocMDP} or \texttt{SigFieldLock} permission level. To prevent SFDA, it should require that no two signature fields reference the same signature dictionary, or, as a weaker but still sufficient alternative, that all fields referencing a given signature dictionary already exist when the dictionary is created.

At the verifier level, individual implementations can adopt both of the above rules without waiting for the specification to change. This precedent already exists: many of the attacks reported in prior work~\cite{pdf-certification-attack,shadow-attacks,processing-dangerous-paths,one-trillion-dollar-refund,practical-decryption-exfiltration} were addressed by individual verifier vendors rather than by amendments to ISO~32000. Verifier-level mitigation, however, does not scale since every new verifier implementation can reintroduce the same class of vulnerability.

At the detection-tooling level, we provide a Python tool that takes a PDF as input and produces a textual report indicating whether either attack is present. When an attack is detected, the report identifies the affected signature fields and the offending PDF objects. ASA is detected by comparing each signature field as it appears in the revision that signed it against its appearance in the latest revision; if \texttt{/AP} or \texttt{/Rect} differ, the document is flagged. SFDA is detected by scanning the latest revision for pairs of signature fields that reference the same signature dictionary and verifying, for each pair, whether at least one of the fields was introduced in an incremental update later than the dictionary it references.

\section{Related Works} \label{sec:related-works}

\cite{one-trillion-dollar-refund} published a large study on approval signatures. They found three attacks that break verifiers. The Universal Signature Forgery (USF) directly changes the signature dictionary so the verifier cannot find what it needs but still reports the signature as valid. The Incremental Saving Attack (ISA) appends a broken incremental update that inserts the visual modification. The Signature Wrapping Attack (SWA) moves the signed bytes elsewhere in the file and puts new content in their place.

\cite{shadow-attacks} builds on that work with the introduction of shadow attacks, which use spec-compliant incremental updates instead of broken ones. Their attack plants hidden content in the document before it is signed and reveals it afterwards through normal incremental updates.

\cite{pdf-certification-attack} looks at certification signatures and finds two attacks. The Evil Annotation Attack (EAA) adds annotations that hide or replace certified content. The Sneaky Signature Attack (SSA) adds a new signature whose appearance covers attacker-controlled content. Both attacks abuse changes that \texttt{DocMDP} allows. SSA also requires the attacker to get a fresh certificate trusted by the verifier, since the new signature has to be considered valid.

ASA differs from these attacks in how it works and in what it asks of the attacker. Unlike USF, ISA, and SWA, ASA uses solely incremental updates fully within the spec. Unlike the shadow attacks, ASA does not need access to the document before signing, and it can be applied to any signed PDF that already has an existing signature. The closest match is EAA, which also covers the page with attacker-chosen content. The difference is that EAA uses a new annotation, which depends on the annotation permissions of \texttt{DocMDP} and \texttt{SigFieldLock}. ASA changes the signature field's own \texttt{/AP} and \texttt{/Rect}, which affected verifiers treat as part of filling the field. This sidesteps the permission checks even at a lower \texttt{DocMDP} or \texttt{SigFieldLock} level than EAA.

SFDA also works only on already-signed PDFs and follows the spec. The closest match is SSA, since both add a new signature field that the verifier shows as valid. The difference is that SSA brings its own signature dictionary, which forces the attacker to get a trusted certificate. SFDA points the new field at a signature dictionary already present in the document, so no extra certificate is needed. This widens the threat to settings where a trusted certificate is hard to get, such as documents signed under closed PKIs. The duplicated field also shows up under the original signer's name and certificate chain, which makes the forgery harder to dismiss than the unfamiliar identity that SSA produces.

\section{Conclusion}

PDF signatures provide a standardized and consistent way for someone to approve the state of a document at a specific instant of time. However, visual modifications can be made by incrementally updating the document in a way that keeps the signature itself valid and triggers no modification warning. In this paper, we introduced two novel attacks that abuse flaws in the specification, the Appearance Substitution Attack (ASA) and the Signature Field Duplication Attack (SFDA). Although neither modifies the page directly, both abuse the appearance mechanism of interactive forms to overlay visual changes on top of the page, and even when detectable the result is subtle enough that most users would not hesitate to trust the signature, since affected verifiers still report it as valid and do not flag the modifications. We then evaluated these attacks against many verifiers and found that most were vulnerable in one way or another. Although we proposed countermeasures to mitigate the damage such attacks can cause, both this work and previous attacks can be seen as proof that the large complexity of the PDF specification, coupled with the decentralized ecosystem where each verifier implements its own algorithm and interprets the ISO in its own way, creates fertile ground for attacks of this kind. In fact, one could argue that incremental updates, which greatly complicate signature validation, are themselves a dangerous functionality, given that both the attacks introduced here and more than half of those in Section~\ref{sec:related-works} abuse incremental updates in one way or another. As future work, it would be valuable to define a standardized verification algorithm on which all verifiers could base their implementations.

After this paper is reviewed, the authors will notify each affected verifier and observe the usual 90-day interval between private disclosure and public release.


\bibliographystyle{sbc}
\bibliography{sbc-template}

\end{document}


